\chapter*{Homework 4}
\section*{Problem 1}
Let $X$ be a random variable with values in $[0,+\infty]$ 
such that $\mathbb{E}[X]=0$. 
Explain why $X<\infty$ almost surely and show that $X=0$ almost surely.

\begin{proof}

Since $X \ge 0$ and $\mathbb{E}[X]=0<\infty$, by Markov's inequality, for any $t>0$,
\[
\mathbb{P}(X \ge t) \le \frac{\mathbb{E}[X]}{t} = 0
\]
Hence $\mathbb{P}(X \ge t)=0$ for all $t>0$. In particular,
\[
\mathbb{P}(X = +\infty) = \lim_{t \to +\infty} \mathbb{P}(X \ge t) =  \lim_{t \to +\infty} 0 = 0
\]
that is, $X<\infty$ almost surely.


Also notice that
\[
\{X>0\} = \bigcup_{n=1}^\infty \left\{X \ge \frac{1}{n}\right\}
\]
and by countable subadditivity,
\[
\mathbb{P}(X>0) \le \sum_{n=1}^\infty \mathbb{P}\left(X \ge \frac{1}{n}\right) = 0
\]
Thus $\mathbb{P}(X>0)=0$.
And since $X$ takes values in $[0,+\infty]$,
 we have $$\bP(\{X=0\}) = \bP(\{X \le 0\}) = 1 - \bP(\{X>0\}) = 1$$, and

This finishes the proof that $X<\infty$ a.s. and $X=0$ a.s.
\end{proof}


\section*{Problem 2}
Let $X$ be a random variable with $\bE[X] = 3$ 
and $\bE[X^2] = 13$. Show that:
$$ 
\bP(-2\leq X \leq 8) \geq \frac{21}{25}
 $$   
\begin{solution}
    
Compute the variance of $X$:
\[
\mathrm{Var}(X) = \mathbb{E}[X^2] - (\mathbb{E}[X])^2 = 13 - 3^2 = 4
\]
Observe that
\[
\mathbb{P}(-2 \le X \le 8) = \mathbb{P}(|X-3| \le 5)
\]
By Chebyshev's inequality,
\[
\mathbb{P}(|X-3| \ge 5) \le \frac{\mathrm{Var}(X)}{5^2} = \frac{4}{25}
\]
Therefore,
\[
\mathbb{P}(|X-3| \le 5) = 1 - \mathbb{P}(|X-3| \ge 5) \ge 1 - \frac{4}{25} = \frac{21}{25}
\]
Thus,
\[
\mathbb{P}(-2 \le X \le 8) \ge \frac{21}{25}
\]
\end{solution}

\section*{Problem 3}
Let $X, Y$ be two random variables such that
 $\mathbb{P}(Y=1)=1 / 5, \mathbb{P}(Y=2)=3 / 5$ 
and $\mathbb{P}(Y=3)= 1 / 5$. In addition
$$
X|\{Y=1\} \sim \operatorname{Exp}(2), X|\{Y=2\} \sim \operatorname{Exp}(3) \text { and } X \mid\{Y=3\}=7 .
$$
Compute the moment generating function of $X$.
\begin{solution}
Compute each conditional moment generating function:

For $X\mid \{Y=1\}\sim \operatorname{Exp}(2)$
\[
\mathbb{E}[e^{tX}\mid Y=1]=\frac{2}{2-t}, \quad t<2
\]
For $X\mid \{Y=2\}\sim \operatorname{Exp}(3)$
\[
\mathbb{E}[e^{tX}\mid Y=2]=\frac{3}{3-t}, \quad t<3
\]
For $X\mid \{Y=3\}=7$, 
\[
\mathbb{E}[e^{tX}\mid Y=3]=e^{7t}
\]
Then we use the law of total expectation conditioning on $Y$: 
for any $t$ such that the expectations below are finite, we have
\begin{align*}
M_X(t)=\mathbb{E}[e^{tX}]
&=\sum_{y=1}^3 \mathbb{E}[e^{tX}\mid Y=y]\mathbb{P}(Y=y)\\
&=\frac{1}{5}\cdot \frac{2}{2-t}
+\frac{3}{5}\cdot \frac{3}{3-t}
+\frac{1}{5}e^{7t}\\
&=\frac{2}{5(2-t)}+\frac{9}{5(3-t)}+\frac{1}{5}e^{7t},
\quad t<2
\end{align*}
So the moment generating function of $X$ is
\[
M_X(t)=\frac{2}{5(2-t)}+\frac{9}{5(3-t)}+\frac{1}{5}e^{7t},
\quad t<2
\]
\end{solution}

\section*{Problem 4}
For any $n \in \mathbb{N}$ with $n \geq 1$ we set $a_n=1 /\left(2 n^2\right)$. 
Consider the sequence of random variables $\left(X_n\right)_{n \in \mathbb{N}}$ with
$$
X_n= 
\begin{cases}
0, & \text { with probability } a_n \\ 
1, & \text { with probability } 1-2 a_n \\ 
n^2, & \text { with probability } a_n
\end{cases}
$$
Check if $\left(X_n\right)$ converges in probability and 
if $\left(X_n\right)$ converges in $L^1$.
\begin{solution}

We first show that $X_n \to 1$ in probability.

Let $\varepsilon>0$.\\
For $n$ large enough s.t. $n^2-1>\varepsilon$,
$|X_n-1|>\varepsilon$ iff $X_n=0$ or $X_n=n^2$. Therefore,
\[
\mathbb{P}(|X_n-1|>\varepsilon)
=\mathbb{P}(X_n=0)+\mathbb{P}(X_n=n^2)
=a_n+a_n
=2a_n
=\frac{1}{n^2}
\]
Since
\(
\frac{1}{n^2}\to 0
\), we have: 
\[
\lim_{n\to\infty} \
\bP(|X_n -1| > \varepsilon) = 0
\]
Since $\varepsilon>0$ is arbitrary, we conclude that
$$
X_n \xrightarrow{\mathbb{P}} 1
$$

We then show that $X_n$ does not converge to $1$ in $L^1$.

We compute
\begin{align*}
\mathbb{E}[|X_n-1|]
&=|0-1| a_n + |1-1| (1-2a_n) + |n^2-1|  a_n\\
&=a_n + (n^2-1)a_n\\
&= n^2 a_n = \frac{1}{2} \not\to 0
\end{align*}
Hence
\[
X_n \not\xrightarrow{L^1} 1
\]
\end{solution}

\section*{Problem 5}
Let $X$ and $Y$ be independent random variables with densities
$$
f_X(x)=\left\{
\begin{array}{ll}
2 x, & 0 \leq x \leq 1, \\
0, & \text { otherwise }
\end{array} \quad f_Y(y)= 
\begin{cases}1 / 2, & 0 \leq y \leq 2 \\
0, & \text { otherwise }
\end{cases}\right.
$$
Find the distribution function of the sum $Z=X+Y$.

\begin{solution}

Since $X$ and $Y$ are independent, the density of
\(
Z=X+Y
\)
is given by convolution:
\[
f_Z(z)=\int_{-\infty}^{\infty} f_X(x)f_Y(z-x)\,dx
\]
where we know 
\[
f_X(x)=2x \mathbf{1}_{[0,1]}(x), \quad f_Y(y)=\frac12 \mathbf{1}_{[0,2]}(y)
\]
Thus
\begin{align*}
f_Z(z)
&=\int_{-\infty}^{\infty} 2x \cdot \frac12 
\mathbf{1}_{[0,1]}(x)\mathbf{1}_{[0,2]}(z-x)\,dx\\
&=\int_{-\infty}^{\infty}x  \bigg(\mathbf{1}_{[0,1]}(x)\bigg)  \bigg( \mathbf{1}_{[0,2]}(z-x)\bigg)   \,dx\\
&=\int_{\max(0,z-2)}^{\min(1,z)} x\,dx
\end{align*}

Compute this piecewise:
If $0\le z\le 1$, then the interval is $[0,z]$, so
\[
f_Z(z)=\int_0^z x\,dx=\frac{z^2}{2}
\]
If $1\le z\le 2$, then the interval is $[0,1]$, so
\[
f_Z(z)=\int_0^1 x\,dx=\frac12
\]
If $2\le z\le 3$, then the interval is $[z-2,1]$, so
\[
f_Z(z)=\int_{z-2}^1 x\,dx
=\frac{1-(z-2)^2}{2}
\]
Outside $[0,3]$, clearly $f_Z(z)=0$.
Therefore,
\[
f_Z(z)=
\begin{cases}
0, & z<0,\\
\dfrac{z^2}{2}, & 0\le z\le 1,\\
\dfrac12, & 1\le z\le 2,\\
\dfrac{1-(z-2)^2}{2}, & 2\le z\le 3,\\
0, & z>3
\end{cases}
\]
Now we integrate to get the cdf.
For $z<0$, $F_Z(z)=0$.

For $0\le z\le 1$,
\[
F_Z(z)=\int_0^z \frac{t^2}{2}\,dt=\frac{z^3}{6}
\]
For $1\le z\le 2$,
\[
F_Z(z)=F_Z(1)+\int_1^z \frac12\,dt
=\frac16+\frac{z-1}{2}
=\frac{z}{2}-\frac13
\]
For $2\le z\le 3$,
\[
F_Z(z)=F_Z(2)+\int_2^z \frac{1-(t-2)^2}{2}\,dt = \frac23+\frac{z-2}{2}-\frac{(z-2)^3}{6}
\]
And for $z\ge 3$, $F_Z(z)=1$.

Hence the cdf of $Z=X+Y$ is
\[
F_Z(z)=
\begin{cases}
0, & z<0,\\
\dfrac{z^3}{6}, & 0\le z\le 1,\\
\dfrac{z}{2}-\dfrac13, & 1\le z\le 2,\\
\dfrac23+\dfrac{z-2}{2}-\dfrac{(z-2)^3}{6}, & 2\le z\le 3,\\
1, & z\ge 3
\end{cases}
\]
\end{solution}


\section*{Problem 6}
Let $a_1, a_2, \ldots, a_n$ and $\lambda$ be positive constants and 
let $\left\{X_i: 1 \leq i \leq n\right\}$ be independent random variables with
$$
X_i \sim \Gamma\left(a_i, \lambda\right), \quad i=1,2, \ldots, n
$$
(i.e., with the same second parameter $\lambda$ ). 
Show that $X_1+X_2+\cdots+X_n \sim \Gamma(a, \lambda)$, 
with $a=\sum_{i=1}^n a_i$.
\begin{proof}

Let
\[
S_n:=X_1+X_2+\cdots+X_n,
\quad a=\sum_{i=1}^n a_i
\]
We need to show that $
S_n \sim \Gamma(a,\lambda)
$.

Since \(X_i \sim \Gamma(a_i,\lambda)\), its moment generating function is
\[
M_{X_i}(t)=\mathbb{E}[e^{tX_i}]
=\left(\frac{\lambda}{\lambda-t}\right)^{a_i},
\quad t<\lambda
\]

Since \(X_1,\dots,X_n\) are independent, 
the moment generating function of \(S_n\) is
\[
M_{S_n}(t)
=\mathbb{E}[e^{t(X_1+\cdots+X_n)}]
=\prod_{i=1}^n \mathbb{E}[e^{tX_i}]
=\prod_{i=1}^n M_{X_i}(t)
\]
Therefore,
\[
M_{S_n}(t)
=\prod_{i=1}^n \left(\frac{\lambda}{\lambda-t}\right)^{a_i}
=\left(\frac{\lambda}{\lambda-t}\right)^{\sum_{i=1}^n a_i}
=\left(\frac{\lambda}{\lambda-t}\right)^a,
\quad t<\lambda
\]

Note this is exactly the moment generating function of a
 \(\Gamma(a,\lambda)\) random variable. Hence,
\[
S_n = 
X_1+\cdots+X_n \sim \Gamma(a,\lambda),
\quad a=\sum_{i=1}^n a_i
\]
\end{proof}